\section*{Hoofdstuk \ref{ch:b3:pericyclic} --- Pericyclische reacties}

\begin{solution}{exo:b3:pericyclic:1}
Diels-Alder: \omterm{def:b3:pericyclic:families}{cycloadditie} [4+2]. Opening van cyclobuteen: elektrocyclisch.
Cope: sigmatroop [3,3]. Sulfoleen dat $\ce{SO2}$ verliest: cheletroop. Ozon en
alkeen: \omterm{def:b3:pericyclic:dipolar}{1,3-dipolaire cycloadditie}
([3+2]). Verschuiving $[1,5]$-H: sigmatroop.
\end{solution}

\begin{solution}{exo:b3:pericyclic:2}
Hexatrieen (6 elektronen): thermisch \omterm{def:b3:pericyclic:rotation-modes}{disrotatorisch}, fotochemisch
\omterm{def:b3:pericyclic:rotation-modes}{conrotatorisch}. Butadieen (4): thermisch \omterm{def:b3:pericyclic:rotation-modes}{conrotatorisch}, fotochemisch
\omterm{def:b3:pericyclic:rotation-modes}{disrotatorisch}.
\end{solution}

\begin{solution}{exo:b3:pericyclic:3}
Vier elektronen, thermisch: conrotatorische opening. De twee
methylgroepen, aan tegenovergestelde vlakken van de ring, draaien beide
naar buiten: (\textit{E},\textit{E})-hexa-2,4-dieen.
\end{solution}

\begin{solution}{exo:b3:pericyclic:4}
De suprafaciale reactie [2+2] heeft $4q$ elektronen: in de grondtoestand overlappen de HOMO van het ene
etheen en de LUMO van het andere met één bindend en één antibindend
uiteinde. In de aangeslagen toestand speelt de enkelvoudig bezette $\pi^*$ van één molecuul de rol van de
HOMO, en passen beide uiteinden.
\end{solution}

\begin{solution}{exo:b3:pericyclic:5}
Het \omterm{def:b3:point-groups:mirror}{spiegelvlak} blijft behouden. Hexatrieen: $\psi_1$ S, $\psi_2$ A, $\psi_3$ S bezet; $\psi_4$ A, $\psi_5$ S, $\psi_6$ A leeg.
Cyclohexadieen, in volgorde: $\sigma$ S, $\pi_1$ S, $\pi_2$ A bezet; $\pi_3^*$ S, $\pi_4^*$ A, $\sigma^*$ A leeg. Verbinden volgens
symmetrie en volgorde: $\psi_1 \to \sigma$, $\psi_2 \to \pi_2$, $\psi_3 \to \pi_1$: bezet naar bezet, de disrotatorische sluiting
is thermisch toegestaan.
\end{solution}

\begin{solution}{exo:b3:pericyclic:6}
Zes elektronen onder licht: \omterm{def:b3:pericyclic:rotation-modes}{conrotatorisch}, zodat de twee eindstandige
methylgroepen aan tegenovergestelde vlakken eindigen:
\textit{trans}-5,6-dimethylcyclohexa-1,3-dieen (de thermische reactie geeft
het isomeer \textit{cis}).
\end{solution}

\begin{solution}{exo:b3:pericyclic:7}
Het waterstofatoom op C5 verhuist naar C1 van het dieen, zijn buur rond
de ring, via een overgangstoestand met zes elektronen ($\sigma2_s + \pi4_s$), \omterm{def:b3:pericyclic:faciality}{suprafaciaal}: één migratie
geeft 1-methyl-, een tweede 2-methylcyclopentadieen. Een keten van vijf
koolstofatomen laat het waterstofatoom aan één vlak blijven, wat de
kleine ring sowieso afdwingt.
\end{solution}

\begin{solution}{exo:b3:pericyclic:8}
$\pi2_s$ (alkeen) $+ \pi2_a$ (C=C van het keteen, aan tegenovergestelde vlakken aangevallen via de
loodrechte $\pi^*$ van C=O): het aantal componenten $(4q + 2)_s$ is één (het alkeen), oneven:
thermisch toegestaan.
\end{solution}

\begin{solution}{exo:b3:pericyclic:9}
In de stoel liggen de C=C van de vinylgroep (atomen 1--2), het zuurstofatoom (3), het koolstofatoom van $\ce{OCH2}$
(4) en de C=C van het allyl (5--6); de binding O--C breekt en de binding C1--C6 vormt zich:
het product is pent-4-enal,
$\ce{OHC-CH2-CH2-CH=CH2}$.
\end{solution}

\begin{solution}{exo:b3:pericyclic:10}
Acht elektronen is $4q$; met één antarafaciale component heeft de
rangschikking een möbiustopologie, aromatisch met $4q$ elektronen: de overgangstoestand is aromatisch en de
reactie thermisch toegestaan (zoals de conrotatorische sluiting van
octatetraeen).
\end{solution}

\begin{solution}{exo:b3:pericyclic:11}
De grootste coëfficiënten koppelen verbindt N3 met het eindstandige
koolstofatoom: het 1,4-digesubstitueerde triazool. Zonder katalysator
verschillen de coëfficiënten weinig en ontstaan zowel het isomeer 1,4 als
1,5; koper(I) laat de reactie stapsgewijs verlopen via een
koperacetylide en geeft alleen het isomeer 1,4.
\end{solution}

\begin{solution}{exo:b3:pericyclic:12}
Het waterstofatoom overbrugt de uiteinden van de SOMO van een radicaal van $j$ atomen, $k = (j + 1)/2$. Zijn
coëfficiënten aan de uiteinden zijn evenredig met $\sin(k\pi/(j + 1))$ en $(-1)^{k+1}\sin(k\pi/(j + 1))$. Voor $j = 7$, $k = 4$: tegengestelde tekens, dus
liggen de lobben met gelijk teken aan de twee uiteinden aan
tegenovergestelde vlakken: \omterm{def:b3:pericyclic:faciality}{antarafaciaal}. Voor
$j = 5$, $k = 3$: hetzelfde teken aan hetzelfde vlak: \omterm{def:b3:pericyclic:faciality}{suprafaciaal}.
\end{solution}

\begin{solution}{pb:b3:pericyclic:1}
\textbf{1.} De $\sigma$-binding C9--C10 van ring B; met het 5,7-dieen geeft ze een hexatrieen,
C10=C5--C6=C7--C8=C9.

\textbf{2.} Zes elektronen (de twee $\pi$-bindingen en de $\sigma$-binding): een elektrocyclische ringopening.

\textbf{3.} Thermisch komen een cyclohexadieen en zijn open hexatrieen
alleen traag in evenwicht, via een hoge barrière, en de ring is de
stabielere kant; het dieen absorbeert ultraviolet licht, en zijn
aangeslagen toestand opent zich binnen picoseconden.

\textbf{4.} \omterm{def:b3:pericyclic:rotation-modes}{Conrotatorisch}: in de aangeslagen toestand heeft de enkelvoudig bezette LUMO van het
trieensysteem ($k = 4$) uiteinden met tegengestelde tekens.

\textbf{5.} De binding C6=C7 komt uit de ring: haar twee ketenvoortzettingen, C5 en C8,
lagen aan dezelfde kant ervan, en blijven daar: \textit{Z}.

\textbf{6.} De thermische opening, alleen \omterm{def:b3:pericyclic:rotation-modes}{disrotatorisch} toegestaan, heeft
een barrière die ver boven wat de lichaamswarmte kan leveren ligt, en zou
bergop leiden, naar het minder stabiele trieen.

\textbf{7.} Van C19 (de methylgroep op C10) naar C9: een sigmatrope verschuiving $[1,7]$-H.

\textbf{8.} Acht: de zes $\pi$-elektronen van het trieen en de twee van de binding C--H.

\textbf{9.} $\sigma2 + \pi6$. \omterm{def:b3:pericyclic:faciality}{Suprafaciaal} zouden beide s zijn, met elk één component $(4q + 2)_s$:
even, verboden. Met het trieen \omterm{def:b3:pericyclic:faciality}{antarafaciaal}, $\sigma2_s + \pi6_a$: één getelde component, oneven,
toegestaan.

\textbf{10.} Zeven atomen in een helixvormige keten met configuratie \textit{Z} laten het waterstofatoom van het ene vlak naar
het andere overgaan; een keten van drie atomen kan niet zo ver
verdraaien.

\textbf{11.} Möbius (één antarafaciale component) met acht elektronen, $4q$: aromatisch.

\textbf{12.} Ze is thermisch toegestaan; haar barrière wordt bij
lichaamstemperatuur traag overschreden.

\textbf{13.} Zes elektronen onder licht: opnieuw \omterm{def:b3:pericyclic:rotation-modes}{conrotatorisch}, maar de
uiteinden kunnen de andere kant op draaien en H9 en de methylgroep C19 aan
de andere vlakken plaatsen: een gesloten stereo-isomeer, lumisterol.

\textbf{14.} Met de centrale binding \textit{E} liggen C19 en C9 aan tegenovergestelde kanten van de keten en ver
uit elkaar: het waterstofatoom kan er niet bij.

\textbf{15.} Previtamine D$_3$ absorbeert zelf licht en wordt omgezet in lumisterol en tachysterol,
die geen vitamine D geven; er wordt een fotostationair mengsel bereikt,
dat de productie aftopt.

\textbf{16.} De opening gebeurt in de aangeslagen toestand, binnen
femtoseconden tot picoseconden; de verschuiving is een thermische reactie met een barrière van \qty{85}{kJ/mol}.

\textbf{17.} $t_{1/2} = \ln 2/k = 0.693/\qty{0.023}{h^{-1}} = \qty{30}{h}$.

\textbf{18.} $1 - \eu^{-0.023 \times 24} = 0.42$.

\textbf{19.} $\ln 10/k = 2.303/\qty{0.023}{h^{-1}} = \qty{100}{h}$.

\textbf{20.} $k(\qty{293.15}{K}) = k(\qty{310.15}{K})\,\eu^{-(E_a/R)(1/293.15 - 1/310.15)} = 0.023 \times \eu^{-1.91} = \qty{3.4e-3}{h^{-1}}$.

\textbf{21.} $2.303/\qty{3.4e-3}{h^{-1}} = \qty{680}{h}$, ongeveer 28 dagen.

\textbf{22.} De thermische stap is bij \qty{37}{\celsius} zeven keer sneller dan bij \qty{20}{\celsius}, zodat het
previtamine dat op een ochtend in de zon gemaakt wordt, binnen enkele
dagen omgezet is.

\textbf{23.} Het steroïdskelet, met zijn vele stereocentra, komt kant en
klaar uit het sterol; licht en warmte voeren daarna de twee
pericyclische stappen uit, precies zoals in de huid.

\textbf{24.} Bij lichaamstemperatuur wordt 90\,\% van het previtamine D$_3$ vitamine D$_3$ in ongeveer
\qty{100}{h}, vier dagen.
\end{solution}
